% =====================================================================
%  Báo cáo Checkpoint 01 - Dự án Web E-commerce (bản tiếng Việt)
%  Biên dịch bằng pdfLaTeX (mặc định của Overleaf). Đổi tên file PDF
%  thành Checkpoint_01_Report_<TenNhom>.pdf trước khi nộp.
%  Sửa các giá trị trong khối "THÔNG TIN NHÓM" bên dưới trước.
% =====================================================================
\documentclass[12pt,a4paper]{article}   % cỡ chữ 12

% --- Font: nên chọn Compiler = XeLaTeX trên Overleaf (Menu > Compiler) ---
\usepackage{iftex}
\ifPDFTeX
  % pdfLaTeX: hiển thị được tiếng Việt nhưng không có Times New Roman
  \usepackage[T5]{fontenc}
  \usepackage[utf8]{inputenc}
\else
  % XeLaTeX/LuaLaTeX: dùng Times New Roman, nếu máy chủ không có thì dùng
  % TeX Gyre Termes (bản sao của Times, nhìn gần như giống hệt)
  \usepackage{fontspec}
  \IfFontExistsTF{Times New Roman}
    {\setmainfont{Times New Roman}}
    {\setmainfont{TeX Gyre Termes}}
\fi
\usepackage[vietnamese]{babel}
\usepackage[a4paper,margin=2.2cm]{geometry}
\usepackage{amssymb}
% --- Thư viện bảng ---
\usepackage{array}             % định dạng cột: >{...}, p{}, m{}, b{}
\usepackage{tabularx}          % bảng tự co giãn theo chiều rộng trang (cột X)
\usepackage{booktabs}          % đường kẻ đẹp: \toprule \midrule \bottomrule \cmidrule
\usepackage{longtable}         % bảng dài tự ngắt sang trang sau
\usepackage{ltablex}           % kết hợp tabularx + longtable (cột X trên bảng nhiều trang)
\keepXColumns                  % giữ chiều rộng cột X như tabularx
\usepackage{multirow}          % gộp ô theo hàng: \multirow{2}{*}{...}
\usepackage{makecell}          % xuống dòng trong ô: \makecell{dòng 1\\dòng 2}
\usepackage[table]{xcolor}     % tô màu hàng/ô: \rowcolor, \cellcolor, \rowcolors
\usepackage{caption}           % chú thích bảng/hình: \caption{...}
\usepackage{float}             % đặt bảng đúng chỗ: \begin{table}[H]
\usepackage{enumitem}
\usepackage{tikz}
\usetikzlibrary{positioning,arrows.meta,shapes.geometric}
\usepackage{hyperref}
\hypersetup{colorlinks=true,linkcolor=blue!50!black,urlcolor=blue!60!black}

\setlist{nosep,leftmargin=1.4em}
\setlength{\parskip}{0.4em}
\setlength{\parindent}{0pt}
\renewcommand{\arraystretch}{1.25}
\newcolumntype{Y}{>{\raggedright\arraybackslash}X}               % cột co giãn, căn trái
\newcolumntype{L}[1]{>{\raggedright\arraybackslash}p{#1}}        % cột cố định, căn trái
\newcolumntype{C}[1]{>{\centering\arraybackslash}p{#1}}          % cột cố định, căn giữa
\newcolumntype{M}[1]{>{\raggedright\arraybackslash}m{#1}}        % cột cố định, căn giữa theo chiều dọc

\newcommand{\yes}{$\checkmark$}
\newcommand{\no}{$\times$}
\newcommand{\half}{$\circ$}

% ------------------------- THÔNG TIN NHÓM (SỬA) ----------------------
\newcommand{\TeamName}{WEB\_2627302\_10}
\newcommand{\ProjectName}{VoltCart}
\newcommand{\RepoURL}{\url{https://github.com/your-team/voltcart}}
\newcommand{\MemberA}{Nguyễn Công Thành\newline (MSSV: ........)}
\newcommand{\MemberB}{Nguyễn Minh Mẫn\newline (MSSV: ........)}
\newcommand{\MemberC}{Lương Nhất Khôi\newline (MSSV: ........)}
\newcommand{\MemberD}{Huỳnh Nhật Hòa\newline (MSSV: ........)}
% ---------------------------------------------------------------------

\begin{document}

% ============================== TIÊU ĐỀ ==============================
\begin{center}
  {\Large\bfseries Checkpoint 01 -- Thành lập nhóm và Nghiên cứu ban đầu}\\[0.4em]
  {\large Dự án Web E-commerce: \textbf{\ProjectName} -- Cửa hàng Điện thoại \& Laptop}\\[0.3em]
  Nhóm: \textbf{\TeamName} \quad|\quad Repository: \RepoURL
\end{center}
\vspace{-0.5em}
\noindent\rule{\linewidth}{0.4pt}

% =====================================================================
\part*{\large Phần A -- Định hướng dự án}
% =====================================================================

\section{Tổng quan dự án}

\begin{tabularx}{\linewidth}{|l|Y|}
\hline
\textbf{Tên dự án}            & \ProjectName \\ \hline
\textbf{Danh mục sản phẩm}    & Thiết bị điện tử tiêu dùng mới: \emph{điện thoại}, \emph{laptop}
                                và một số \emph{phụ kiện} cơ bản (sạc, tai nghe, ốp lưng). \\ \hline
\textbf{Khách hàng mục tiêu}  & Sinh viên và người mới đi làm (18--35 tuổi) tại Việt Nam, cần mua
                                điện thoại hoặc laptop để học tập và làm việc, thường so sánh cấu hình
                                trên mạng và muốn đặt hàng nhanh mà không bắt buộc đăng ký tài khoản. \\ \hline
\end{tabularx}

\textbf{Mô tả vấn đề.}
Các chuỗi bán lẻ điện tử lớn hiển thị hàng trăm mẫu máy với bảng thông số dài, kèm nhiều
banner trả góp và khuyến mãi, khiến việc mua một sản phẩm đơn giản trở nên chậm và khó hiểu
với người mua lần đầu. Nhiều khách hàng cũng bỏ ngang khi bị yêu cầu tạo tài khoản.
\ProjectName{} hướng đến một cửa hàng tập trung, nơi khách hàng có thể lọc thiết bị theo
vài thông số quan trọng nhất (hãng, giá, RAM, bộ nhớ), thấy một mức giá rõ ràng và đặt hàng
với tư cách khách vãng lai chỉ trong vài bước, trong khi thành viên đã đăng ký được hưởng giá
thành viên thấp hơn.

\textbf{Các quyết định về phạm vi:}
\begin{itemize}
  \item \textbf{Không làm biến thể sản phẩm.} Mỗi cấu hình là một sản phẩm riêng
        (ví dụ ``Laptop X 16GB/512GB'' và ``Laptop X 8GB/256GB'' là hai sản phẩm,
        mỗi sản phẩm có tồn kho riêng).
  \item \textbf{Thông số chính cố định}, lưu thành cột (hãng, RAM, bộ nhớ, kích thước màn hình)
        và một bảng thông số dạng văn bản để hiển thị ở trang chi tiết.
  \item \textbf{Quy tắc giá:} mỗi sản phẩm có \emph{giá khách} và \emph{giá thành viên}
        (giá thành viên $\le$ giá khách). Thành viên đã đăng nhập luôn thấy và thanh toán theo giá thành viên.
  \item \textbf{Một quy tắc coupon:} giảm số tiền cố định cho đơn đạt giá trị tối thiểu
        (ví dụ \texttt{VOLT200K}: giảm 200{.}000đ cho đơn từ 3{.}000{.}000đ), có trạng thái bật/tắt.
  \item \textbf{Phí giao hàng cố định} 30{.}000đ và \textbf{thanh toán giả lập}
        (thanh toán khi nhận hàng / chuyển khoản giả lập).
  \item \textbf{Tổng tiền} = Tạm tính (giá áp dụng $\times$ số lượng) $-$ Giảm giá coupon $+$ Phí giao hàng.
        Giá được lưu vào đơn hàng tại thời điểm đặt để việc đổi giá sau này không ảnh hưởng đơn cũ.
  \item \textbf{Ngoài phạm vi:} trả góp, thu cũ đổi mới, cổng thanh toán thật, quản lý IMEI/bảo hành,
        nhiều kho hàng, nhiều vai trò quản trị.
\end{itemize}

\section{Thành viên và phân công}

Mỗi thành viên phụ trách trọn vẹn một mảng để ai cũng có thể giải thích và demo phần việc của mình.

\begin{tabularx}{\linewidth}{|L{3.3cm}|L{2.7cm}|Y|}
\hline
\rowcolor{gray!15} \textbf{Thành viên} & \textbf{Vai trò} & \textbf{Nhiệm vụ chính} \\ \hline \endhead
\MemberA & Trưởng nhóm, Backend lõi & Thiết kế CSDL, xác thực và phân quyền, logic giá (khách/thành viên),
          giỏ hàng và thanh toán, kiểm tra coupon, quy trình Git và review code. \\ \hline
\MemberB & Frontend khách hàng & Viết HTML/CSS thuần: layout, hệ thống giao diện (biến CSS dùng chung),
          responsive, trang chủ, danh sách sản phẩm có
          tìm kiếm/lọc/sắp xếp, trang chi tiết, trang giỏ hàng cập nhật bằng AJAX. \\ \hline
\MemberC & Khu vực quản trị & Layout admin, CRUD sản phẩm và danh mục kèm ảnh, quản lý giá và tồn kho,
          quản lý đơn hàng và cập nhật trạng thái, danh sách khách hàng, quản lý coupon. \\ \hline
\MemberD & Báo cáo, tài khoản, kiểm thử & Dashboard và báo cáo doanh thu, hồ sơ khách hàng và lịch sử đơn,
          dữ liệu mẫu/seeder, \texttt{TEST\_CASES.xlsx}, tài liệu, video demo. \\ \hline
\end{tabularx}

\section{Danh sách tính năng ban đầu}

\begin{tabularx}{\linewidth}{|L{3.2cm}|Y|}
\hline
\rowcolor{gray!15} \textbf{Nhóm chức năng} & \textbf{Tính năng} \\ \hline \endhead
Duyệt sản phẩm & Trang chủ có banner và sản phẩm nổi bật; danh sách sản phẩm theo danh mục
           (Điện thoại, Laptop, Phụ kiện); tìm kiếm theo từ khóa; lọc theo hãng, khoảng giá, RAM;
           sắp xếp theo giá/mới nhất; trang chi tiết có ảnh, thông số, giá khách và giá thành viên,
           tình trạng còn hàng. \\ \hline
Tài khoản  & Đăng ký, đăng nhập, đăng xuất có kiểm tra dữ liệu; hồ sơ (họ tên, email, số điện thoại,
           địa chỉ); lịch sử đơn hàng và chi tiết đơn kèm trạng thái. \\ \hline
Giỏ hàng \& thanh toán & Thêm vào giỏ, cập nhật số lượng, xóa sản phẩm (AJAX); kiểm tra tồn kho;
           form thanh toán cho khách vãng lai và thành viên; áp dụng coupon (AJAX); tóm tắt giá (tạm tính,
           giảm giá, phí giao hàng, tổng cộng); thanh toán giả lập; trang xác nhận đơn hàng. \\ \hline
Quản trị viên & Khu vực admin được bảo vệ; dashboard (số đơn, doanh thu, khách hàng, sản phẩm, doanh thu
           theo ngày); CRUD sản phẩm/danh mục kèm ảnh, giá, tồn kho, trạng thái; cảnh báo khi xóa sản phẩm
           đã có trong đơn hàng; tìm kiếm đơn, xem chi tiết, cập nhật trạng thái (Chờ xử lý, Đã xác nhận,
           Đang giao, Hoàn thành, Đã hủy); danh sách khách hàng và trạng thái tài khoản; quản lý coupon;
           báo cáo doanh thu theo khoảng thời gian. \\ \hline
Nâng cao (chỉ làm khi phần cốt lõi đã ổn định) & Khách vãng lai tra cứu đơn bằng mã đơn + email;
           báo cáo sản phẩm sắp hết hàng; so sánh hai thiết bị (mục tiêu mở rộng). \\ \hline
\end{tabularx}

\section{Luồng người dùng và Wireframe}

\begin{center}
\begin{tikzpicture}[
  node distance=0.45cm and 0.45cm,
  box/.style={draw,rounded corners=2pt,align=center,font=\footnotesize,
              minimum height=0.85cm,text width=1.75cm,fill=blue!6},
  adm/.style={box,fill=orange!10},
  dec/.style={draw,diamond,aspect=2,align=center,font=\footnotesize,inner sep=1pt,fill=yellow!10},
  arr/.style={-{Stealth[length=2mm]},thick}
]
% luồng khách hàng
\node[box] (home) {Trang chủ};
\node[box,right=of home] (list) {Danh sách\\tìm/lọc};
\node[box,right=of list] (detail) {Chi tiết sản phẩm};
\node[box,right=of detail] (cart) {Giỏ hàng\\(AJAX, coupon)};
\node[dec,right=of cart] (logged) {Đã đăng\\nhập?};
\node[box,below=0.9cm of logged] (guest) {Form khách\\giá khách};
\node[box,below=0.9cm of cart] (member) {Form thành viên\\giá TV};
\node[box,left=of member] (pay) {Thanh toán giả lập};
\node[box,left=of pay] (confirm) {Xác nhận\\đơn hàng};
\node[box,left=of confirm] (status) {Lịch sử / trạng thái đơn};

\draw[arr] (home) -- (list);
\draw[arr] (list) -- (detail);
\draw[arr] (detail) -- (cart);
\draw[arr] (cart) -- (logged);
\draw[arr] (logged) -- node[right,font=\scriptsize]{không} (guest);
\draw[arr] (logged.south west) -- node[above left,font=\scriptsize]{có} (member.north east);
\draw[arr] (guest) -- (member);
\draw[arr] (member) -- (pay);
\draw[arr] (pay) -- (confirm);
\draw[arr] (confirm) -- (status);

% luồng quản trị
\node[adm,below=1.3cm of status] (alogin) {Đăng nhập admin};
\node[adm,right=of alogin] (dash) {Dashboard};
\node[adm,right=of dash] (prod) {Sản phẩm \&\\danh mục};
\node[adm,right=of prod] (ord) {Đơn hàng: cập nhật trạng thái};
\node[adm,right=of ord] (rep) {Coupon \&\\báo cáo};
\draw[arr] (alogin) -- (dash);
\draw[arr] (dash) -- (prod);
\draw[arr] (prod) -- (ord);
\draw[arr] (ord) -- (rep);
\draw[arr,dashed] (ord.north) -- ++(0,0.45) -| node[pos=0.25,above,font=\scriptsize]{khách thấy trạng thái mới} (status.south);
\end{tikzpicture}

{\small Hình 1. Luồng mua hàng của khách (xanh; nhánh khách vãng lai đi qua form khách) và luồng quản trị (cam).}
\end{center}

\begin{center}
\begin{tikzpicture}[font=\scriptsize]
  \draw (0,0) rectangle (12,6.2);
  % header
  \draw[fill=gray!15] (0,5.5) rectangle (12,6.2);
  \node[anchor=west] at (0.2,5.85) {\textbf{\ProjectName}};
  \draw (3.2,5.65) rectangle (8.2,6.05); \node at (5.7,5.85) {Tìm điện thoại, laptop...};
  \node at (9.3,5.85) {Đăng nhập};
  \node at (10.8,5.85) {Giỏ hàng (2)};
  % bộ lọc
  \draw[fill=gray!5] (0.2,0.3) rectangle (2.8,5.3);
  \node[anchor=north west,align=left] at (0.3,5.2)
       {\textbf{Bộ lọc}\\[2pt] Danh mục\\ $\square$ Điện thoại\\ $\square$ Laptop\\
        $\square$ Phụ kiện\\[3pt] Hãng\\ $\square$ Hãng A\\ $\square$ Hãng B\\[3pt]
        Khoảng giá\\ {[}\quad--\quad{]}\\[3pt] RAM\\ $\square$ 8GB \ $\square$ 16GB};
  % thanh sắp xếp
  \node[anchor=west] at (3.1,5.05) {24 sản phẩm};
  \node[anchor=east] at (11.8,5.05) {Sắp xếp: Giá thấp--cao $\blacktriangledown$};
  % lưới sản phẩm
  \foreach \x in {0,1,2}{
    \foreach \y in {0,1}{
      \begin{scope}[shift={(3.1+\x*2.95,0.3+\y*2.3)}]
        \draw (0,0) rectangle (2.7,2.15);
        \draw[fill=gray!20] (0.15,0.95) rectangle (2.55,2.0);
        \node at (1.35,1.47) {ảnh};
        \node at (1.35,0.75) {Tên sản phẩm};
        \node[font=\tiny] at (1.35,0.45) {Giá khách: 12.990.000đ};
        \node[font=\tiny] at (1.35,0.18) {\textbf{Giá TV: 12.490.000đ}};
      \end{scope}}}
\end{tikzpicture}

{\small Hình 2. Wireframe trang danh sách sản phẩm (máy tính). Trên điện thoại, cột bộ lọc thu gọn
thành nút ``Bộ lọc'' và lưới sản phẩm chuyển thành một hoặc hai cột.}
\end{center}

\section{Công nghệ đề xuất}

\begin{tabularx}{\linewidth}{|l|Y|}
\hline
\rowcolor{gray!15} \textbf{Thành phần} & \textbf{Lựa chọn và lý do} \\ \hline \endhead
Backend   & PHP 8.2 + Laravel 11 (môn học khuyến nghị; có sẵn xác thực, validation, migration, seeder). \\ \hline
Frontend  & HTML5 + CSS3 tự viết, không dùng framework CSS: Flexbox/Grid cho bố cục, biến CSS
            (\texttt{:root}) cho màu và font thống nhất, media query cho responsive. Giao diện được đặt
            trong Blade template của Laravel để hiển thị dữ liệu từ CSDL. JavaScript thuần (\texttt{fetch})
            cho AJAX (số lượng trong giỏ, coupon, gợi ý tìm kiếm). \\ \hline
Cơ sở dữ liệu & MySQL 8 (chạy local bằng XAMPP/Laragon); lược đồ quản lý bằng migration của Laravel. \\ \hline
Biểu đồ   & Chart.js cho biểu đồ doanh thu ở dashboard (không bắt buộc với phần cốt lõi). \\ \hline
Công cụ   & Git + GitHub (nhánh theo tính năng, pull request), Figma hoặc draw.io cho wireframe, VS Code. \\ \hline
\end{tabularx}

\section{Kế hoạch phát triển (8 tuần)}

\begin{tabularx}{\linewidth}{|l|Y|}
\hline
\rowcolor{gray!15} \textbf{Tuần} & \textbf{Công việc chính} \\ \hline \endhead
1--2 & Lập nhóm, chọn danh mục sản phẩm, nghiên cứu website tham khảo, xây dựng persona, vẽ wireframe,
       tạo repository Git theo cấu trúc thư mục bắt buộc. \\ \hline
3 (CP1) & Trình bày Checkpoint 01; chốt thiết kế CSDL (ERD); khởi tạo project Laravel; viết file CSS
       chung (màu, font, nút, form, lưới) và layout HTML cơ bản (header, footer, menu). \\ \hline
4    & Migration và seeder với dữ liệu mẫu thực tế; xác thực; danh sách và chi tiết sản phẩm lấy từ CSDL;
       admin CRUD sản phẩm và danh mục. \\ \hline
5    & Giỏ hàng bằng AJAX; thanh toán cho khách vãng lai; giá thành viên; quản lý đơn hàng cơ bản ở admin. \\ \hline
6 (CP2) & Demo luồng cốt lõi; quy tắc coupon; cập nhật trạng thái đơn và hiển thị cho khách;
       hoàn thiện responsive; ghi chú tiến độ. \\ \hline
7    & Dashboard và báo cáo doanh thu; hồ sơ và lịch sử đơn; thông báo lỗi/validation;
       viết và chạy test case; bắt đầu tính năng nâng cao đã chọn. \\ \hline
8 (CP3) & Dừng thêm tính năng; sửa lỗi; hoàn thiện \texttt{TEST\_CASES.xlsx}, báo cáo cuối kỳ,
       hướng dẫn cài đặt và video demo; nộp bài. \\ \hline
\end{tabularx}

\section{Repository và công việc ban đầu}

Repository: \RepoURL. Repository đã tuân theo cấu trúc bắt buộc
(\texttt{CODE/}, \texttt{DATABASE/}, \texttt{DATA/}, \texttt{DOCUMENTATION/}, \texttt{VIDEO/},
\texttt{README.md}). Quy trình làm việc: nhánh \texttt{main} (ổn định) và \texttt{develop},
mỗi công việc một nhánh tính năng (ví dụ \texttt{feature/cart-ajax}), gộp qua pull request
có thành viên khác review. Các commit ban đầu: cấu trúc thư mục, README, \texttt{STUDENT\_INFO.md},
wireframe và báo cáo này.

\newpage
% =====================================================================
\part*{\large Phần B -- Báo cáo nghiên cứu (Mini-Research)}
% =====================================================================

\section{Website tham khảo}

Nhóm đã khảo sát ba nhà bán lẻ điện tử phổ biến tại Việt Nam, đều bán cả điện thoại và laptop:

\begin{enumerate}
  \item \textbf{Thế Giới Di Động} -- \url{https://www.thegioididong.com}:
        chuỗi bán lẻ toàn quốc với danh mục sản phẩm rất lớn; bộ lọc mạnh và bảng thông số
        chi tiết ở mỗi trang sản phẩm.
  \item \textbf{FPT Shop} -- \url{https://fptshop.com.vn}:
        nhà bán lẻ lớn với bố cục sản phẩm gọn gàng, nhiều khuyến mãi và tính năng so sánh sản phẩm.
  \item \textbf{CellphoneS} -- \url{https://cellphones.com.vn}:
        được sinh viên ưa chuộng; có chương trình thành viên (Smember) với giá ưu đãi riêng,
        rất gần với yêu cầu giá khách/giá thành viên của dự án.
\end{enumerate}

\section{Khách hàng mục tiêu và Persona}

Khách hàng mục tiêu là sinh viên và người mới đi làm mua thiết bị để học tập hoặc làm việc,
có ngân sách hạn chế và chủ yếu mua sắm trên điện thoại.

\begin{tabularx}{\linewidth}{|L{2.4cm}|Y|Y|}
\hline
\rowcolor{gray!15}  & \textbf{Persona 1 -- Minh (khách vãng lai)} & \textbf{Persona 2 -- Lan (thành viên)} \\ \hline \endhead
Hồ sơ       & 19 tuổi, sinh viên CNTT năm nhất, ngân sách khoảng 15 triệu đồng. &
              27 tuổi, nhân viên văn phòng, hay mua phụ kiện và đổi điện thoại mỗi 2--3 năm. \\ \hline
Mục tiêu    & Tìm laptop có ít nhất 16GB RAM để học lập trình. &
              Mua được giá tốt nhất và theo dõi đơn hàng dễ dàng. \\ \hline
Khó khăn    & Quá nhiều mẫu máy và banner; không muốn tạo tài khoản chỉ để mua một lần. &
              Phải nhập lại địa chỉ mỗi lần mua; không rõ giá cuối cùng sau khuyến mãi. \\ \hline
Cần ở \ProjectName{} & Bộ lọc đơn giản (giá, RAM), tình trạng hàng rõ ràng, thanh toán nhanh không cần tài khoản. &
              Giá thành viên, lưu hồ sơ và địa chỉ, lịch sử và trạng thái đơn hàng. \\ \hline
\end{tabularx}

\section{So sánh tính năng}

% TODO (nhóm): kiểm tra lại từng ô trên website thật và chụp màn hình đưa vào phụ lục.
\begin{center}
\begin{tabularx}{\linewidth}{|Y|c|c|c|c|}
\hline
\rowcolor{gray!15} \textbf{Tính năng} & \textbf{TGDĐ} & \textbf{FPT Shop} & \textbf{CellphoneS} & \textbf{\ProjectName} \\ \hline \endhead
Menu danh mục và sản phẩm nổi bật          & \yes & \yes & \yes & \yes \\ \hline
Tìm kiếm có gợi ý                          & \yes & \yes & \yes & \yes \\ \hline
Lọc theo hãng / giá / thông số             & \yes & \yes & \yes & \yes \\ \hline
Sắp xếp theo giá                           & \yes & \yes & \yes & \yes \\ \hline
Bảng thông số chi tiết                     & \yes & \yes & \yes & \yes \\ \hline
Thông tin tồn kho / còn hàng               & \yes & \yes & \yes & \yes \\ \hline
Đặt hàng không cần tài khoản               & \yes & \yes & \yes & \yes \\ \hline
Giá thành viên / chương trình thành viên   & \half & \half & \yes & \yes \\ \hline
Mã giảm giá / voucher                      & \yes & \yes & \yes & \yes (1 quy tắc) \\ \hline
Theo dõi đơn hàng                          & \yes & \yes & \yes & \yes \\ \hline
So sánh sản phẩm                           & \yes & \yes & \yes & mở rộng \\ \hline
Trả góp                                    & \yes & \yes & \yes & \no \\ \hline
Thu cũ đổi mới                             & \yes & \yes & \yes & \no \\ \hline
Đánh giá và chấm điểm sản phẩm             & \yes & \yes & \yes & \no \\ \hline
Cổng thanh toán trực tuyến thật            & \yes & \yes & \yes & \no{} (giả lập) \\ \hline
\end{tabularx}

{\small \yes{} có \quad \half{} một phần hoặc không nổi bật \quad \no{} không làm}
\end{center}

\section{Tính năng được chọn và lý do}

\textbf{Được chọn.}
\begin{itemize}
  \item \emph{Lọc theo hãng, giá và RAM}: cả ba website đều dựa nhiều vào bộ lọc, và Persona 1
        quyết định chủ yếu theo ngân sách và RAM. Nhóm chỉ giữ vài bộ lọc để đơn giản.
  \item \emph{Hiển thị rõ giá khách và giá thành viên ở thẻ sản phẩm, trang chi tiết, giỏ hàng và
        thanh toán}: lấy cảm hứng từ giá Smember của CellphoneS; tạo lý do để khách đăng ký
        (Persona 2) và đáp ứng yêu cầu của môn học.
  \item \emph{Thanh toán không cần tài khoản}: cả ba website đều hỗ trợ; giải quyết khó khăn chính của Persona 1.
  \item \emph{Giỏ hàng và coupon bằng AJAX}: phản hồi ngay khi đổi số lượng và tổng tiền mà không tải lại trang.
  \item \emph{Lịch sử và trạng thái đơn hàng}: cần thiết để thành viên theo dõi việc giao hàng.
\end{itemize}

\textbf{Không chọn (và lý do).}
Trả góp, thu cũ đổi mới, cổng thanh toán thật và nhiều loại khuyến mãi rất phổ biến ở các nhà bán
lẻ lớn nhưng cần dịch vụ bên ngoài và quy tắc phức tạp, không phù hợp với dự án 8 tuần. Biến thể
sản phẩm được thay bằng các sản phẩm riêng cho từng cấu hình. Đánh giá sản phẩm và so sánh sản phẩm
chỉ là phần nâng cao tùy chọn, làm sau khi luồng cốt lõi đã ổn định.

\section{Nhận xét UI/UX}

\begin{itemize}
  \item \textbf{Điểm mạnh nên học theo:} header cố định với thanh tìm kiếm lớn; thẻ sản phẩm hiển thị
        ảnh, tên, thông số chính và giá trong lưới đồng nhất; nhãn màu cho tình trạng hàng;
        bộ lọc thu gọn thành nút trên điện thoại; trang thanh toán ngắn, luôn hiển thị tóm tắt đơn hàng.
  \item \textbf{Điểm yếu cần tránh:} quá nhiều banner khuyến mãi và pop-up đẩy sản phẩm xuống dưới;
        nhiều mức giá trên một thẻ (giá cũ, giá online, giá trả góp) gây rối cho người mua;
        trang sản phẩm quá dài trên điện thoại.
  \item \textbf{Định hướng thiết kế của nhóm:} bố cục gọn với một màu nhấn chủ đạo, tối đa hai mức giá
        mỗi sản phẩm (giá khách và giá thành viên, ghi nhãn rõ ràng), có thông báo khi đang tải,
        thành công, không có kết quả và lỗi; lưới responsive ưu tiên thiết bị di động. Toàn bộ giao diện
        do nhóm tự viết bằng HTML/CSS (không dùng template có sẵn) để đảm bảo tính nguyên bản.
\end{itemize}

\textbf{Link tham khảo:}
\begin{itemize}
  \item Thế Giới Di Động: \url{https://www.thegioididong.com}
  \item FPT Shop: \url{https://fptshop.com.vn}
  \item CellphoneS: \url{https://cellphones.com.vn}
\end{itemize}

\end{document}
